\documentclass[a4paper,11pt]{article}
\usepackage{fontspec}
\usepackage{polyglossia}
\setmainlanguage{russian}
\setmainfont{Times New Roman}
\newfontfamily\cyrillicfont{Times New Roman}
\usepackage[margin=2.2cm]{geometry}
\usepackage{setspace}
\usepackage{longtable,array,booktabs}
\setstretch{1.1}
\setlength{\parindent}{0pt}
\setlength{\parskip}{4pt}

\begin{document}

\begin{center}
{\Large Календарный план научно-исследовательской работы}\\[4pt]
{\large на период 13 ноября – 12 декабря 2026 г.}\\[8pt]
Кадыров Э.Э., группа 420, кафедра ИИТ, ВМК МГУ\\
Научные руководители: <научные руководители>\\[4pt]
Тема: «Прогноз космологического красного смещения по фотометрическим признакам»
\end{center}

\begin{longtable}{>{\raggedright\arraybackslash}p{2.9cm}>{\raggedright\arraybackslash}p{7.2cm}>{\raggedright\arraybackslash}p{5.2cm}}
\toprule
Сроки & Задача & Результат \\
\midrule
\endhead
13–19 ноя (1 нед) &
Обновление презентации ВКР, часть 1: постановка (обучение на полях F/S, тест на поле N), данные и метки, перевод фотометрии F/S в полосы поля N, схема модели (предобучение CLIP и дообучение). &
Обновлённые слайды постановки, данных и модели \\
20–26 ноя (1 нед) &
Обновление презентации ВКР, часть 2: результаты на поле N (DESI, NEP-Wide, HETDEX) в сравнении с Phosphoros и PHZ Euclid, точность по интервалам $z$ и яркости, слабые места, проверенные варианты модели; согласование с научным руководителем. &
Обновлённая презентация со ссылкой в репозитории \\
27 ноя – 1 дек (5 дней) &
Опыты с CLIP: новые архитектуры энкодеров в мультимодальной сети. &
Метрики на тестовых наборах поля N в сравнении с итоговой моделью \\
2–5 дек (4 дня) &
Опыты с CLIP: варианты предобучения (данные, модальности, длительность). &
Та же таблица метрик для новых вариантов \\
6–10 дек (5 дней) &
Сравнение лучших вариантов с итоговой моделью, Phosphoros и PHZ Euclid на поле N. &
Сводная таблица; решение, какая версия становится основной \\
11–12 дек (2 дня) &
Сведение метрик и отчётов, отчёт за период, план на следующий месяц. &
Сводная таблица, отчёт, план \\
\bottomrule
\end{longtable}

\end{document}
